\section{Background}

\subsection{Linear Attention}

Given $Q, K, V \in \mathbb{R}^{B \times H \times L \times d_h}$, linear 
attention computes:

\begin{equation}
O_t = \frac{\phi(q_t)^\top S_t}{\phi(q_t)^\top z_t}
\end{equation}

where:

\begin{equation}
S_t = \sum_{j \leq t} \phi(k_j) v_j^\top, \quad
z_t = \sum_{j \leq t} \phi(k_j)
\end{equation}

$S_t \in \mathbb{R}^{F \times d_h}$ and $z_t \in \mathbb{R}^F$ are the 
running state and normalizer.

\subsection{Exponential Decay Recurrence}

We add a per-token, per-head decay $\lambda_t \in (0, 1)$:

\begin{equation}
S_t = \lambda_t S_{t-1} + \phi(k_t) v_t^\top
\end{equation}

\begin{equation}
z_t = \lambda_t z_{t-1} + \phi(k_t)
\end{equation}

Unrolled:

\begin{equation}
S_t = \sum_{j \leq t} \left( \prod_{i=j+1}^{t} \lambda_i \right) \phi(k_j) v_j^\top
\end{equation}

The coefficient of contribution from position $j$ to position $t$ is:

\begin{equation}
D_{t, j} = \prod_{i=j+1}^{t} \lambda_i = \exp\!\left(\sum_{i=j+1}^{t} \log \lambda_i\right)
\end{equation}

This is an \textbf{exponential decay} with rate controlled by $\lambda$.

\subsection{Chunked Scan Algorithm}

For a chunk of $C$ tokens, we carry state $S_\text{prev}$ from the previous 
chunk. The output within the chunk is a combination of:

\begin{itemize}
\item \textbf{Intra-chunk contributions:} dense causal attention within the chunk.
\item \textbf{Inter-chunk contributions:} attention to the state $S_\text{prev}$.

\end{itemize}
The intra-chunk scores are:

\begin{equation}
\text{scores}_{ij}^\text{intra} = q_i^\top k_j \cdot 
\exp\!\left(\sum_{p=j+1}^{i} \log \lambda_p\right), \quad j \leq i
\end{equation}

The inter-chunk scores:

\begin{equation}
\text{scores}_{i}^\text{inter} = q_i^\top S_\text{prev} \cdot 
\exp\!\left(\sum_{p=1}^{i} \log \lambda_p\right)
\end{equation}

Then $S_\text{new}$ is computed by decaying $S_\text{prev}$ by the total 
chunk decay and adding the intra-chunk contributions.

\subsection{ELU+1 Feature Map}

\begin{equation}
\phi(x) = \text{elu}(x) + 1 = \begin{cases}
x + 1 & \text{if } x > 0 \\
e^x & \text{if } x \leq 0
\end{cases}
\end{equation}

\textbf{Property.} $\phi(x) > 0$ for all $x$, which ensures $z_t > 0$ and the 
normalizer is well-behaved.

\subsection{Data-Dependent Decay}

We learn $\lambda_{h,t}$ from the input:

\begin{equation}
\lambda_{h,t} = \sigma(W_g x_t + b_g + \ell_h)
\end{equation}

where:

\begin{itemize}
\item $W_g \in \mathbb{R}^{H \times d}$, $b_g \in \mathbb{R}^H$ -- learned gate.
\item $\ell_h \in \mathbb{R}$ -- per-head bias, initialized to linearly spaced
  values.

\end{itemize}
\textbf{Initialization.} $W_g = 0$, $b_g = 0$, $\ell_h = \text{linspace}(-0.5, 0.0, H)$.

\textbf{Effect.} At step 0, $\lambda_{h,t} = \sigma(\ell_h)$ -- constant per head, 
giving each head a fixed "default horizon".

\subsection{Bidirectional Attention}

For encoders, we run both forward and backward scans:

\begin{equation}
O_t = \frac{1}{2} \left( O_t^\text{fwd} + O_t^\text{bwd} \right)
\end{equation}

Each $O_t^\text{fwd}$, $O_t^\text{bwd}$ uses the same recurrence but in 
opposite directions.

\subsection{Grouped-Query Attention (GQA)}

For GQA, we expand KV heads before the scan:

\begin{equation}
K, V \in \mathbb{R}^{B \times H_{kv} \times L \times d_h}, \quad
K_\text{expanded} = \text{repeat\_kv}(K, H / H_{kv})
\end{equation}

\subsection{Notation}

\begin{table}[H]
\centering
\small
\begin{tabular}{ll}
\toprule
\textbf{Symbol} & \textbf{Meaning} \\
\midrule
$\lambda_{h,t}$ & Decay at position $t$, head $h$ \\
$\ell_h$ & Per-head bias \\
$W_g, b_g$ & Gate parameters \\
$S, z$ & Running state \\
$\phi$ & ELU+1 feature map \\
$C$ & Chunk size \\
\bottomrule
\end{tabular}
\end{table}

